\section*{第 \ref{ch:b2:heart} 章——心脏与心动周期}

\begin{solution}{exo:b2:heart:1}
腔\omterm{def:b2:blood-circulation:vessels}{静脉} $\to$ 右\omterm{def:b2:heart:anatomy}{心房} $\to$ 三尖瓣 $\to$ 右
\omterm{def:b2:heart:anatomy}{心室} $\to$ 肺\omterm{def:b2:blood-circulation:vessels}{动脉}瓣 $\to$ 肺\omterm{def:b2:blood-circulation:vessels}{动脉} $\to$ 肺
$\to$ 肺\omterm{def:b2:blood-circulation:vessels}{静脉} $\to$ 左\omterm{def:b2:heart:anatomy}{心房} $\to$ 二尖瓣 $\to$ 左
\omterm{def:b2:heart:anatomy}{心室} $\to$ 主\omterm{def:b2:blood-circulation:vessels}{动脉}瓣 $\to$ 主\omterm{def:b2:blood-circulation:vessels}{动脉}。
\end{solution}

\begin{solution}{exo:b2:heart:2}
充盈期（\omterm{prop:b2:heart:cycle}{舒张期}）：房室瓣开放，半月瓣关闭。等容
\omterm{prop:b2:heart:cycle}{收缩期}：四个瓣膜全部关闭。射血期：半月瓣开放，房室瓣关闭。
等容\omterm{prop:b2:heart:cycle}{舒张期}：全部关闭。第一心音为\omterm{prop:b2:heart:cycle}{收缩期}开始时房室瓣
关闭所致；第二心音为半月瓣
在\omterm{prop:b2:heart:cycle}{收缩期}末关闭所致。
\end{solution}

\begin{solution}{exo:b2:heart:3}
\omterm{prop:b2:heart:conduction}{窦房结}（$t = 0$；\omterm{def:b2:heart:anatomy}{心房}去极化时为 P 波）；
\omterm{prop:b2:heart:conduction}{房室结}（约在 \qty{80}{ms} 时到达，冲动在此
停留 \qty{100}{ms}，即 PR 段）；希氏束及
\omterm{prop:b2:heart:conduction}{浦肯野纤维}（\qtyrange{180}{220}{ms}；\omterm{def:b2:heart:anatomy}{心室}去极化时为 QRS 波群）；
随后\omterm{def:b2:heart:anatomy}{心室}复极化时出现 T 波。
\end{solution}

\begin{solution}{exo:b2:heart:4}
在工作范围内，\omterm{def:b2:heart:anatomy}{心室}充盈得越多，射出的血量
也越多。若右\omterm{def:b2:heart:anatomy}{心室}泵出更多\omterm{def:b2:blood-circulation:blood}{血液}，左\omterm{def:b2:heart:anatomy}{心室}
几搏之后便接收到更多\omterm{def:b2:blood-circulation:blood}{血液}，并依同一定律泵出更多。任何
失衡都会在无信号的情况下自行纠正。
\end{solution}

\begin{solution}{exo:b2:heart:5}
\omterm{prop:b2:heart:cycle}{每搏输出量} \qty{80}{mL}；\omterm{prop:b2:heart:cycle}{射血分数} $80/130 = \qty{62}{\%}$；
\omterm{thm:b2:heart:starling}{心输出量} $80\times 70 = \qty{5.6}{L/min}$。衰竭时：\qty{60}{mL}，
$60/180 = \qty{33}{\%}$，\qty{4.2}{L/min}——扩张的\omterm{def:b2:heart:anatomy}{心室}
以较大容积工作，使\omterm{thm:b2:heart:starling}{心输出量}基本维持，但射出的血量
仅为其所容纳血量的三分之一。
\end{solution}

\begin{solution}{exo:b2:heart:6}
$W = 100\times 133\times 7\times 10^{-5} = \qty{0.93}{J}$；功率
$0.93\times 1.2 = \qty{1.1}{W}$。在 \qty{150}{mmHg} 时：\qty{1.4}{J}，
\qty{1.7}{W}。额外的 \qty{0.56}{W} 机械功率对应
\qty{3.7}{W} 代谢功率，每秒耗氧 \qty{0.19}{mL}，即额外 \qty{11}{mL/min}。
\end{solution}

\begin{solution}{exo:b2:heart:7}
$60/0.5 = 120$ 次/分；$60/1.5 = 40$。第三种为完全性\omterm{def:b2:heart:anatomy}{心脏}
传导阻滞——\omterm{def:b2:heart:anatomy}{心房}在\omterm{prop:b2:heart:conduction}{窦房结}控制下以 100 次/分搏动，\omterm{def:b2:heart:anatomy}{心室}
以 33 次/分搏动，由各自的\omterm{prop:b2:heart:conduction}{起搏点}驱动，\omterm{prop:b2:heart:conduction}{房室结}不传导任何冲动。
\end{solution}

\begin{solution}{exo:b2:heart:8}
漂移时间 $20/25 = \qty{0.8}{s}$，加上 \qty{0.15}{s}：\qty{0.95}{s}，
即每分钟 63 次。乙酰胆碱：$25/18 = \qty{1.39}{s}$，加上 $0.15$：
\qty{1.54}{s}，每分钟 39 次。去甲肾上腺素：$20/40 = \qty{0.5}{s}$，
加上 $0.15$：\qty{0.65}{s}，每分钟 92 次。
\end{solution}

\begin{solution}{exo:b2:heart:9}
延迟使\omterm{def:b2:heart:anatomy}{心房}在\omterm{def:b2:heart:anatomy}{心室}收缩之前完成向\omterm{def:b2:heart:anatomy}{心室}排血；
若无此延迟，\omterm{def:b2:heart:anatomy}{心房}与\omterm{def:b2:heart:anatomy}{心室}将同时收缩，房室瓣
会在\omterm{def:b2:heart:anatomy}{心室}半充盈时即关闭，\omterm{prop:b2:heart:cycle}{每搏输出量}将因
\omterm{prop:b2:heart:cycle}{每搏输出量}将减去\omterm{def:b2:heart:anatomy}{心房}的
贡献。长 PR 间期只是推迟\omterm{def:b2:heart:anatomy}{心室}\omterm{prop:b2:heart:cycle}{收缩期}，
在延迟增长到使搏动被漏掉之前并无代价，
直至传导完全失效。
\end{solution}

\begin{solution}{exo:b2:heart:10}
静息：$5000/45 = \qty{111}{mL}$；运动：$\num{30000}/180 =
\qty{167}{mL}$。斯塔林定律将运动时较大的\omterm{def:b2:blood-circulation:vessels}{静脉}回流
转化为较大的\omterm{prop:b2:heart:cycle}{每搏输出量}；交感神经则增加
心率与收缩力（使收缩末期容积减小）。约
200 次/分以上时，\omterm{prop:b2:heart:cycle}{舒张期}过短，不足以充满\omterm{def:b2:heart:anatomy}{心室}，\omterm{prop:b2:heart:cycle}{每搏输出量}
下降速度快于心率的上升。
\end{solution}

\begin{solution}{exo:b2:heart:11}
$v = Q/A$：$250/3 = \qty{83}{cm/s}$ 与 $250/1 = \qty{250}{cm/s}$。
$\Delta P = \tfrac12\times 1060\times v^{2}$：\qty{365}{Pa}
（\qty{2.7}{mmHg}）与 \qty{3300}{Pa}（\qty{25}{mmHg}）。\omterm{def:b2:heart:anatomy}{心室}通过增厚心壁
产生额外的压强；增厚而僵硬的\omterm{def:b2:heart:anatomy}{心室}
最终超出其冠状\omterm{def:b2:blood-circulation:vessels}{动脉}的供血能力，充盈
不良，并最终衰竭。
\end{solution}

\begin{solution}{exo:b2:heart:12}
自律性使\omterm{def:b2:heart:anatomy}{心脏}无需任何输入即可搏动；斯塔林定律
使输出量与回流相匹配，无需任何输入即可平衡两侧\omterm{def:b2:heart:anatomy}{心室}；
神经与肾上腺素仅在这一自我调节的核心基础上
设定心率与收缩力。移植\omterm{def:b2:heart:anatomy}{心脏}仍能搏动
（约 100 次/分，无迷走神经张力），并依充盈程度
调节\omterm{prop:b2:heart:cycle}{每搏输出量}；运动时可借斯塔林机制及
循环中的肾上腺素提高输出量，但反应较慢，
且弱于有神经支配的\omterm{def:b2:heart:anatomy}{心脏}。
\end{solution}

\begin{solution}{pb:b2:heart:1}
\textbf{1.} \qty{70}{mL}；$70/130 = \qty{54}{\%}$；$70\times 70 =
\qty{4.9}{L/min}$。
\textbf{2.} $60/70 = \qty{0.86}{s}$；\omterm{prop:b2:heart:cycle}{舒张期} $0.86 - 0.3 =
\qty{0.56}{s}$。
\textbf{3.} $70/0.3 = \qty{233}{mL/s}$；$233/3 = \qty{78}{cm/s}$。
\textbf{4.} $70/1500 = \qty{0.047}{s}$，约 \qty{50}{ms}。
\textbf{5.} $100\times 133\times 7\times 10^{-5} = \qty{0.93}{J}$；
$0.93\times 70/60 = \qty{1.1}{W}$。
\textbf{6.} 两\omterm{def:b2:heart:anatomy}{心室}合计 $1.1\times 1.2 = \qty{1.3}{W}$；按
\qty{15}{\%}计，代谢功率 \qty{8.7}{W}；$8.7/20 = \qty{0.43}{mL}$ 氧/秒，
即 \qty{26}{mL/min}。
\textbf{7.} 每毫升\omterm{def:b2:blood-circulation:blood}{血液}提取 \qty{0.14}{mL} 氧，
$26/0.14 = \qty{190}{mL/min}$。
\textbf{8.} 间隔 \qty{0.86}{s}，漂移时间 \qty{0.71}{s}：
$20/0.71 = \qty{28}{mV/s}$。
\textbf{9.} 间隔 \qty{0.6}{s}，漂移时间 \qty{0.45}{s}：\qty{44}{mV/s}。
\textbf{10.} 间隔 \qty{0.33}{s}，漂移时间 \qty{0.18}{s}：\qty{110}{mV/s}。
\textbf{11.} R--R \qty{0.86}{s}；P：\omterm{def:b2:heart:anatomy}{心房}去极化；QRS：
\omterm{def:b2:heart:anatomy}{心室}去极化；T：\omterm{def:b2:heart:anatomy}{心室}复极化。
\textbf{12.} 延迟使\omterm{def:b2:heart:anatomy}{心房}在\omterm{def:b2:heart:anatomy}{心室}收缩前将\omterm{def:b2:blood-circulation:blood}{血液}排入\omterm{def:b2:heart:anatomy}{心室}；
心率为 180 次/分时整个周期仅 \qty{0.33}{s}，因此
若延迟仍为 \qty{0.16}{s}，就会占去一半——交感
神经同样加快\omterm{prop:b2:heart:conduction}{房室结}的传导。
\textbf{13.} 完全性\omterm{def:b2:heart:anatomy}{心脏}传导阻滞：\omterm{def:b2:heart:anatomy}{心室}以 40 次/分
搏动，由希氏束或\omterm{prop:b2:heart:conduction}{浦肯野纤维}起搏；\omterm{def:b2:heart:anatomy}{心房}在\omterm{prop:b2:heart:conduction}{窦房结}
控制下以 75 次/分搏动。
\textbf{14.} \qty{120}{mL}；$120/150 = \qty{80}{\%}$；$120\times 180 =
\qty{21.6}{L/min}$。
\textbf{15.} $0.33 - 0.3 = \qty{0.03}{s}$：几乎没有充盈时间；
若再加快，\omterm{prop:b2:heart:cycle}{每搏输出量}即骤降。
\textbf{16.} $W = 120\times 133\times 1.2\times 10^{-4} = \qty{1.9}{J}$，
$\times 3$ 每秒：\qty{5.7}{W}；两\omterm{def:b2:heart:anatomy}{心室}合计 \qty{6.9}{W}；
代谢功率 \qty{46}{W}；耗氧 \qty{2.3}{mL/s}，\qty{140}{mL/min}。
\textbf{17.} $140/0.14 = \qty{1000}{mL/min}$，为静息血流的五倍，
须在缩短至周期十分之一的\omterm{prop:b2:heart:cycle}{舒张期}内输送：
冠状\omterm{def:b2:blood-circulation:vessels}{动脉}\omterm{def:b2:blood-circulation:vessels}{小动脉}必须扩张至其极限。
\textbf{18.} 无迷走神经张力时心率约 100 次/分，收缩末期容积
保持 \qty{60}{mL} 不变：\omterm{prop:b2:heart:cycle}{每搏输出量} $150 - 60 = \qty{90}{mL}$，
\omterm{thm:b2:heart:starling}{心输出量} \qty{9}{L/min}。
\textbf{19.} 仅心率变化（$70 \to 180$，每搏量 \qty{70}{mL}，即
\qty{12.6}{L/min}）：$+7.7$；充盈（180 时多 \qty{20}{mL}）：
$+3.6$；收缩性（180 时残余量少 \qty{30}{mL}）：$+5.4$——
由 4.9 增至 \qty{21.6}{L/min}，三项贡献合计
恰好为 $+16.7$。
\textbf{20.} $233/1 = \qty{233}{cm/s}$，即 \qty{2.3}{m/s}。
\textbf{21.} $\tfrac12\times 1060\times 2.33^{2} = \qty{2900}{Pa}$，
即 \qty{22}{mmHg}。
\textbf{22.} $120/0.3 = \qty{400}{mL/s}$，\qty{4}{m/s}；$\tfrac12\times
1060\times 16 = \qty{8500}{Pa}$，\qty{64}{mmHg}。
\textbf{23.} 静息时峰值为 $120 + 22 = \qty{142}{mmHg}$，运动时约为 $140 +
64 = \qty{204}{mmHg}$；静息时平均射血压强由 100 升至 122：每搏功
$\times 1.22$。
\textbf{24.} 较厚的心壁能产生更大的力，并降低每根纤维
所受的应力；然而需灌注的质量增长，而
冠脉血流在\omterm{prop:b2:heart:cycle}{收缩期}受厚壁挤压、\omterm{prop:b2:heart:cycle}{舒张期}又缩短，
不足以跟上需求，厚壁舒张不良，
充盈减少，随之而来的是缺血与衰竭。
\textbf{25.} 静息时 \qty{4.9}{L/min}、运动时 \qty{21.6}{L/min}；
每搏 \qty{0.93}{J}；压强梯度静息时 \qty{22}{mmHg}，运动时
\qty{64}{mmHg}。
\end{solution}
